% titel: Blatt 8, A8: Magnesiumoxid-Elementarzelle (Dichte, Gitterenergie)
% skript: Mortimer Kap. 12, Prüfungsbeilage (Elementarzellen)

% Teil I: Werkzeuge, kein Lösungsweg.

Aus einer Elementarzelle lassen sich Strukturtyp, Stöchiometrie,
Koordinationszahlen und die Dichte ableiten.

\begin{werkzeugkasten}[Elementarzelle auswerten]
\begin{itemize}
  \item \textbf{Strukturtyp} an der Anordnung erkennen (hier das übliche
        Gitter vom Natriumchlorid- bzw.\ Steinsalz-Typ: kubisch-flächen-
        zentiert mit zugehörigen Oktaederlücken).
  \item \textbf{Zahl der Formeleinheiten $Z$} pro Zelle: Teilchen in den
        Ecken (gehören zu $1/8$ in die Zelle), Kanten ($1/4$), Flächen
        ($1/2$), im Inneren (`1`).
  \item \textbf{Koordinationszahl (KZ):} Zahl der nächsten Gegenionen um ein
        Ion --- im Steinsalz-Gitter ist sie für alle Ionen gleich.
  \item \textbf{Dichte} über Zellvolumen $V = a^3$:
        \[
          \rho = \frac{Z \cdot M}{N_A \cdot a^3}.
        \]
  \item \textbf{Gitterenergie:} ~skaliert mit $\frac{q_1 q_2}{r}$; doppelte
        Ladungen und kleine Ionenradien erhöhen sie stark.
\end{itemize}
\textit{Skript:} Mortimer Kap. 12, Prüfungsbeilage (Elementarzellen).
\end{werkzeugkasten}

\begin{achtung}[Darauf achten]
\begin{itemize}
  \item Den Wurffaktor beachten: ein Ecken-Ion gehört nur zu $1/8$, ein
        Flächen-Ion zu $1/2$ in die betrachtete Zelle.
  \item Einheit der Kantenlänge: $a=4.21\cdot10^{-10}\,$m in cm umrechnen.
  \item $M(\mathrm{MgO}) \approx 40.30\ g\,mol^{-1}$; $Z=4$.
  \item Die Koordinationszahl ist räumlich (Struktur), nicht an der
        Zeichnung allein ablesbar.
\end{itemize}
\end{achtung}
